\documentclass[12pt]{article}
\usepackage[utf8]{inputenc}
\usepackage[T1]{fontenc}
\usepackage{hyperref}
\title{Lote Language Resources for K-12 Curriculum Development}
\author{Compiled by Claude Sonnet 5 (claude-sonnet-5) for Andr\'as Kornai}
\date{}
\begin{document}
\maketitle

\section{Introduction}
Lote (ISO 639-3 uvl; also Lohote, formerly Uvol) is an Austronesian language of the Ngero--Vitiaz branch of Western Oceanic. About 6,000 people speak it around Cape Dampier on the south coast of New Britain, Papua New Guinea. No dedicated dictionary turned up in the searches.

\section{Primary Educational Resources}
No primers or school materials were found.

\section{Linguistic and Reference Materials}
\begin{itemize}
\item Wikipedia overview, which cites the main reference: Pearson and van den Berg (2008), \emph{Lote Grammar Sketch}, SIL-PNG Data Papers on Papua New Guinea Languages 54: \url{https://en.wikipedia.org/wiki/Lote_language}
\item R. Barlow's 2024 Lote wordlist, from elicitation with a speaker from Uvol village, as cited in the search summary (archive DOI 10.5281/zenodo.14505730; this DOI link did not itself appear in the search results).
\item ABVD (Austronesian Basic Vocabulary Database) entry: \url{https://abvd.eva.mpg.de/austronesian/language.php?id=1709}
\item Numerals page, Max Planck Institute: \url{https://lingweb.eva.mpg.de/channumerals/Lote.htm}
\end{itemize}

\section{Religious and Community Resources}
The ABVD entry cites a 2009 Wycliffe New Testament translation in Lote, but no direct link to the text was found. The Joshua Project profile is a missions listing: \url{https://joshuaproject.net/index.php/people_groups/15754}

\section{Digital and Audio Resources}
Omniglot page: \url{https://omniglot.com/writing/lote.htm}. The ABVD entry also mentions recordings made by M. Ross in 1977--1982, but no online access was found.

\section{Curriculum Development Context}
Lote is relatively vital, and all community members use it as a first language. It is not known to be taught in schools. A curriculum could draw on the grammar sketch, the wordlist, and the New Testament translation, but no teaching materials turned up and community partnership would be needed.
\end{document}